% NUMBERED PROOF-CHAIN SUMMARY::MAIN14::RP-DIAGNOSTICS::2026-10-21
% Source snapshot: b3ea46290587fcaea559552ed899f942f956b327
\documentclass[11pt]{article}
\usepackage[T1]{fontenc}
\usepackage[utf8]{inputenc}
\usepackage{lmodern,amsmath,amssymb,amsthm,booktabs,enumitem,xurl,hyperref}
\usepackage[margin=.82in]{geometry}
\hypersetup{hidelinks}
\setlength{\parindent}{0pt}
\setlength{\parskip}{5pt}
\setlength{\emergencystretch}{3em}
\setlist{nosep,leftmargin=*}
\newtheorem{theorem}{Theorem}[section]
\newtheorem{proposition}[theorem]{Proposition}
\newcommand{\Pclass}{\mathsf P}
\newcommand{\RP}{\mathsf{RP}}
\newcommand{\NP}{\mathsf{NP}}
\newcommand{\PSPACE}{\mathsf{PSPACE}}
\newcommand{\poly}{\operatorname{poly}}
\newcommand{\SAT}{\operatorname{SAT}}
\newcommand{\repo}[2]{\href{https://github.com/j15marti/pnp/blob/b3ea46290587fcaea559552ed899f942f956b327/#1}{#2}}
\title{\vspace{-1.2cm}Proof-Chain Summary 14\\\large RP wrappers, hitting lists, and the language-level boundary}
\author{P versus NP / P versus RP Research Project}
\date{21 October 2026}
\begin{document}
\maketitle
\vspace{-1em}

\section{Scope and the two unfinished endpoints}
\textbf{P versus RP remains OPEN. P versus NP remains OPEN.}
This document organizes the verified RP diagnostics in reports
RP-HITTING-QUANTIFIER-FIREWALL36, RP-RELATIVIZATION-BIDIRECTIONAL-BARRIER37,
RP-DENSITY-HITTING-TRADEOFF38, and RP-LINEAR-HITTING-TIGHTNESS39
\cite{r36,r37,r38,r39}, with their relevant wrapper ancestors.
It is a proof-chain summary, not a claim of new research or a separation.

\textbf{Source snapshot.} All project links below are pinned to repository
\texttt{j15marti/pnp}, commit
\begin{center}\small\texttt{b3ea46290587fcaea559552ed899f942f956b327}.\end{center}
The statements labelled \textbf{PROVED/EXACT} below have self-contained
 elementary proofs supplied here, restating or clarifying the cited project
results. \textbf{KNOWN THEOREM} identifies external literature.
\textbf{COMPUTATIONAL} means finite evidence only; historical audit counts
are source-reported, not newly rerun. \textbf{CONDITIONAL}, \textbf{OPEN}, and
\textbf{FALSIFIED/RETIRED} retain their literal scope.

\subsection*{The established chain and its stopping points}
\begin{enumerate}
\item A uniform polynomial-time machine with no accepting NO tapes and at
least half accepting YES tapes recognizes an RP language.
\item A fixed such verifier has a polynomial-size length-wise hitting list:
\[
\forall M\ \forall n\ \exists H_{M,n}\ \forall x\in\{0,1\}^n.
\]
This is \emph{nonuniform existence}, not an efficient list constructor.
\item Logarithmic \emph{total} randomness is uniformly enumerable.
A sufficiently efficient machine-specific hitting-list constructor would also
 give deterministic simulation. That general constructor is \textbf{OPEN}.
\item Public-list failure and large hitting requirements for a chosen verifier
can occur even for $\{0,1\}^*$. They do not establish a hard language.
\item Opposite oracle worlds block conclusions from arguments that remain
valid unchanged relative to every oracle. Hitting advice exists in both worlds.
\item Low-degree arithmetic NONIDENTITY is a genuine RP language. The proposed
CNF clause-product route to it fails, and its required deterministic lower bound
is \textbf{OPEN}.
\end{enumerate}

The separation endpoint requires \emph{one fixed uniform language}
$L\in\RP$ such that no deterministic polynomial-time machine decides $L$.
The derandomization endpoint requires a deterministic polynomial-time decider
for every $L\in\RP$. A lower bound for one implementation, a fixed polynomial
clock, a finite portfolio, or a seed-replay interface is neither endpoint.
The basic inclusions
\[
\Pclass\subseteq\RP\subseteq\NP
\]
follow by ignoring randomness and, respectively, using an accepting tape as
an NP witness. Thus $\Pclass\ne\RP$ would imply $\Pclass\ne\NP$.

\newpage
\section{Randomness and wrapper interfaces}
A verifier $M$ is fixed, uniform, and worst-case polynomial-time on every
input and every random tape. Pad its independent fair bits to a declared
polynomial bound $r(n)$, and put
\[
 A_x=\{\rho\in\{0,1\}^{r(n)}:M(x;\rho)=1\}.
\]
The RP contract is $A_x=\varnothing$ for $x\notin L$ and
$|A_x|\ge 2^{r(n)-1}$ for $x\in L$. Construction, sampling, replay, and
interface simulation are charged; random bits are not a free oracle.
Gill's probabilistic-machine model is the classical antecedent \cite{gill}.

\begin{proposition}[Logarithmic-coin simulation: PROVED/EXACT]
If $r(n)=O(\log(n+2))$, exhaustive tape replay and OR give an exact uniform
deterministic polynomial-time decider.
\end{proposition}
\begin{proof}
There are $2^{r(n)}=n^{O(1)}$ tapes. A NO input accepts on none; a YES input
accepts on at least one. Each simulation and all tape handling are polynomial.
\end{proof}
Adaptive completion-time reads are included in the \emph{total} bit budget.
A logarithmic number of bits per stage over polynomially many stages is not
logarithmic total randomness. Amplification must likewise pay every repeated
coin and simulation \cite{seed}.

\begin{proposition}[Fixed-tape zero-error lemma: PROVED/EXACT, restricted]
Suppose a standard fair-stream machine is, for each input $x$, almost surely
correct and almost surely bounded by input-size-only time $B_T(|x|)$ and
simultaneous space $B_S(|x|)$. Each finite execution prefix reads finitely many
stream bits. Then every fixed tape, including the all-zero tape, is correct
and satisfies those bounds, with its interface simulation charged.
\end{proposition}
\begin{proof}
An incorrect output or a first time/space overrun has a finite witnessing
execution prefix. The finitely many observed bits define a cylinder of
positive probability on which the same violation occurs, contradicting the
almost-sure guarantee. Nontermination violates the fixed time bound after a
finite prefix as well.
\end{proof}
This is CLOCKFIX1 \cite{clockfix}, \emph{not} a derandomization of RP or
expected-time Las Vegas algorithms. Waiting for the first $1$ has finite
expectation but diverges on the zero tape. Bounded-error machines can have
incorrect fixed tapes.

\subsection*{Completion-time freshness is not static hardness}
CRS-PPC and COMPLETIONCUT1 \cite{crs,completion} assume that a newly exposed
$\ell$-bit block is uniform conditional on the causal past. For $M$ causal
programs, each fixing at most $q$ guesses before its block is revealed,
\[
 \Pr[\text{some program guesses its block}]\le Mq2^{-\ell}.
\]
Conditioning on each past gives $q2^{-\ell}$; a union bound needs no
independence between the programs. If $Mq<2^\ell$, one common stream defeats
all these \emph{prepredictions}. It need not give one common static SAT input:
solver-dependent completion times can select different later formulas.
A compact post-release function such as $r\mapsto r$ is not a preprediction.
Exact replay needs the actual sampled blocks, and measured latencies when
those affect the controller. These interface and common-input restrictions
cannot be dropped when translating a wrapper result to ordinary decision time.

\newpage
\section{Hitting advice and the missing uniform constructor}
\begin{theorem}[Density-sensitive hitting advice: PROVED/EXACT]
Fix $M,n$, with $m$ YES inputs at length $n$. Suppose every YES acceptance set
has density at least $\alpha>0$ and every NO acceptance set is empty. If an
integer $t\ge1$ satisfies
\[
 m(1-\alpha)^t<1,                                      \tag{1}
\]
then a list of $t$ tapes hits every YES acceptance set. OR replay with this
advice decides the length-$n$ slice exactly.
\end{theorem}
\begin{proof}
Choose $t$ independent uniform tapes. A particular YES row is missed with
probability at most $(1-\alpha)^t$. The union bound over its $m$ YES rows
is strictly below one, so some list misses none. Empty NO rows cannot create
false positives when replayed.
\end{proof}
For rational $\alpha=a/q$, (1) is exactly $m(q-a)^t<q^t$.
For $0<\alpha<1$ and $m\ge1$, the sufficient choice is
\[
 t=\left\lfloor\frac{\log m}{-\log(1-\alpha)}\right\rfloor+1.
\]
If $m=0$, use the empty list; if $\alpha=1$ and $m>0$, one tape suffices.
Literal list storage is $tr(n)$ bits, and replay takes
$O(t(T_M(n)+r(n)))$ time, or $O(tT_M(n))$ when $T_M$ already charges coin
usage. The list can depend on both $M$ and $n$.

At RP density $1/2$, $m\le2^n$ makes $t=n+1$ sufficient. Hence
$\RP\subseteq\Pclass/\mathrm{poly}$, the one-sided advice conclusion
associated with Adleman's theorem \cite{adleman,r36,r38}. Inverse-polynomial
density also gives polynomial advice when $r,T_M$ are polynomial. For the
extreme parameters $m=2^n$ and $\alpha=2^{-n}$, the sufficient certificate
has order $n2^n$. This is not a minimum-list lower bound. Separately, one
random assignment on a unique-witness $n$-variable SAT formula succeeds with
probability $2^{-n}$ and fails the RP density contract.

\subsection*{A precise search interface}
For fixed polynomial-clock $M$, unary $1^n$, and polynomial-size $H$,
invalidity has the certificate
\[
 \exists x,\rho\quad M(x;\rho)=1\ \wedge\
                      \bigwedge_{h\in H}M(x;h)=0.       \tag{2}
\]
All simulations in (2) are polynomial; the conjunction is an explicit
polynomial loop. Thus validity is in coNP. Asking whether a prefix extends
to some valid polynomial-length list is in $\Sigma_2^{\Pclass}$.
These are upper bounds, not hardness results. If the machine code varies,
its simulation requires a suitable explicit unary clock or a fixed common
polynomial exponent; an arbitrary binary-coded exponent is not such a bound.

\textbf{CONDITIONAL.} Under $\Pclass=\NP$, PH collapses and prefix search
constructs the list bit by bit. This only rederives
$\Pclass\subseteq\RP\subseteq\NP=\Pclass$.
An unconditional polynomial-time hitting-list constructor for every fixed RP
verifier would give $\Pclass=\RP$; existence by averaging does not.
Likewise, $\SAT\in\RP$ implies $\NP=\RP$ via deterministic reductions,
not an established deterministic simulation or a conclusion $\Pclass=\RP$.

\newpage
\section{Two trivial languages that defeat overstrong inferences}
\begin{theorem}[Public-list diagonalization: PROVED/EXACT]
Let a uniform polynomial-time procedure produce a polynomial-size list
$H_n\subseteq\{0,1\}^n$. There is an RP verifier for $\{0,1\}^*$ on which
replaying every seed in $H_n$ rejects, for every sufficiently large $n$.
\end{theorem}
\begin{proof}
Eventually $|H_n|\le2^{n-1}$. Define $M_H(x;\rho)$ to reject exactly when
$\rho\in H_{|x|}$. Computing the list and testing membership take polynomial
time. Acceptance has probability at least $1/2$, but every listed seed fails.
On the finitely many exceptional lengths, define $M_H$ to accept every tape.
\end{proof}
The verifier depends on the proposed public-list generator. The conclusion
is that no such machine-independent list works for \emph{all} verifiers;
it does not contradict machine-indexed advice. The language is decided by
one always-accept state. Thus the diagonalization supplies no language outside
P \cite{seed,r36}.

\begin{theorem}[Linear hitting tightness: PROVED/EXACT]
For $n\ge1$ and $x,r\in\mathbb F_2^n$, define
\[
 M(x;r)=\begin{cases}
 1,&x=0,\\
 \langle x,r\rangle,&x\ne0.
 \end{cases}                                         \tag{3}
\]
Its language is $\{0,1\}^*$, its nonzero-input acceptance density is exactly
$1/2$, and a seed set $H$ hits every length-$n$ input iff
$\operatorname{span}(H)=\mathbb F_2^n$. Consequently its minimum seed count
is exactly $n$.
\end{theorem}
\begin{proof}
A nonzero linear functional over $\mathbb F_2^n$ takes each value equally
often. If $\operatorname{span}(H)$ is proper, its orthogonal complement
contains a nonzero $x$, rejected on every $h\in H$. Conversely, if $H$ spans,
no nonzero $x$ is orthogonal to every $h$. The zero input accepts every tape,
and a spanning set is nonempty. A basis has $n$ elements; fewer cannot span.
\end{proof}
This is the rank calculation of RP-LINEAR-HITTING-TIGHTNESS39 \cite{r39}.
The earlier $n+1$ bound is therefore tight to within one seed as a general
\emph{chosen-verifier} bound.

\textbf{Representation qualification.} An explicit list of $n$ seeds of
length $n$ occupies $n^2$ bits. This is not an information-theoretic advice
lower bound: the standard basis is uniformly generable without advice.
Nor is it a time lower bound for the language, which again has an
always-accept decider.

\subsection*{The missing implication}
\[
\boxed{\text{difficulty of seed replay for }M
       \quad\not\Rightarrow\quad
       \text{difficulty of every decider for }L(M).}
\]
A valid language lower bound must exclude alternate algorithms and
representations, including deterministic bypasses. The same warning applies
to proposed hitting lower bounds for a Schwartz--Zippel verifier, to fixed
hashes and prime provenance, and to wrapper-specific transcript obligations.
Report 39's finite rank/replay checks are \textbf{COMPUTATIONAL}; the
linear-algebra argument above, not the finite sample, proves the all-$n$ claim.

\newpage
\section{Two opposite oracle worlds}
\begin{theorem}[Relativization diagnostic: PROVED/EXACT]
There are recursive oracles $A,B$ for which
$\Pclass^A\ne\RP^A$ and $\Pclass^B=\RP^B$.
\end{theorem}
\begin{proof}
Enumerate all deterministic polynomial-time oracle machines with explicit
clocks as $(D_i,T_i)$. At stage $i$, choose $n_i$ larger than the previous
chosen length and every earlier fixed query length, with
$T_i(n_i)<2^{n_i-1}$. Simulate $D_i(1^{n_i})$, respecting all previously
assigned answers and answering every newly queried string by zero forever.

If $D_i$ accepts, leave the whole length-$n_i$ slice empty. If it rejects,
set exactly $2^{n_i-1}$ unqueried strings of that length to one. Fewer than
half the slice were queried, so enough strings remain. Later stages preserve
this slice and every query answer on the simulated computation. All other
slices default to empty.

Define
\[
 L^A=\{1^n:n\ge1,\ |A\cap\{0,1\}^n|=2^{n-1}\}.
\]
Reject nonunary inputs. On $1^n$, sample one uniform $n$-bit string and query
$A$. Every slice is empty or exactly half full, giving exact NO soundness and
YES probability $1/2$. Thus $L^A\in\RP^A$. At $1^{n_i}$, $D_i^A$ gives the
opposite answer, so $L^A\notin\Pclass^A$. The oracle is recursive: to decide
a string of length $n$, simulate the finite stages through the first chosen
length exceeding $n$; no later stage can add a one at length $n$.

For equality, let $B=\mathrm{TQBF}$. PSPACE-completeness and polynomial-space
simulation of polynomially many PSPACE oracle calls give
\[
 \Pclass^B=\PSPACE,\qquad
 \Pclass^B\subseteq\RP^B\subseteq\NP^B\subseteq\PSPACE.
\]
Hence $\Pclass^B=\RP^B$.
\end{proof}
This makes the construction in report 37 explicit \cite{r37}. It is a
method barrier, not an unrelativized separation: no proof valid unchanged
relative to every oracle can establish either equality or inequality.

\subsection*{Advice survives both worlds}
For any fixed oracle $O$, every $\RP^O$ acceptance matrix still has empty
NO rows and half-dense YES rows. Section 3 therefore proves
\[
                 \RP^O\subseteq\Pclass^O/\mathrm{poly}
                 \qquad\text{for every }O.
\]
In particular, the separating oracle has polynomial hitting advice but no
uniform polynomial-time deterministic simulation. Advice existence alone
cannot close that gap.

\subsection*{Primary-literature attribution}
Baker--Gill--Solovay establish opposite oracle worlds for P versus NP
\cite{bgs}; the specific one-sided construction above is proved here, not
inferred from a P-versus-NP separation. Bennett--Gill's random-oracle theorem
actually has $\Pclass^O=\mathsf{BPP}^O$ with probability one \cite{bg}, hence
$\Pclass^O=\RP^O$ there. It must not be cited as a random-oracle separation
of P from RP. A black-box timing or freshness argument that survives all
oracles supplies no missing nonrelativizing step.

\newpage
\section{Arithmetic NONIDENTITY and the CNF failure}
\textbf{KNOWN THEOREM.} For a nonzero polynomial $f$ of total degree at most
$D$ over $\mathbb F_q$, uniform $y\in\mathbb F_q^v$ satisfies
\[
             \Pr[f(y)=0]\le D/q.                      \tag{4}
\]
This is the Schwartz--Zippel bound; see Schwartz and the explicit finite-field
statement and proof of Moshkovitz \cite{schwartz,moshkovitz}.

\begin{proposition}[An explicit RP member: PROVED/EXACT using (4)]
Let an input of length $N$ encode a division-free arithmetic circuit over an
input prime $q\le N^3$, with syntactic degree $D\le N^2$ and
$q\ge8\max(1,D)$. The language of well-formed circuits computing a nonzero
formal polynomial is in RP.
\end{proposition}
\begin{proof}
Reject malformed inputs. Primality can be tested by trial division because
$q\le N^3$; circuit validation, degree bounds and field evaluation have
polynomial bit cost. For each of the $v=O(N)$ occurring variables set
\[
 w=\lceil\log_2q\rceil,\qquad
 K=\lceil\log_2(8\max(1,v))\rceil.
\]
Draw $K$ independent $w$-bit candidates and take the first below $q$; reject
if none exists. Each coordinate fails with probability at most $2^{-K}$,
so all succeed with probability at least $7/8$. Conditional on success the
coordinates are independent uniform residues. Equation (4) then gives
nonzero evaluation with probability at least $7/8$ on a YES input. Overall
acceptance is at least $49/64$. Zero polynomials never accept.
The total budget is exactly $vKw=O(N\log^2N)$ fair bits, with unused samples
ignored; constant circuits are evaluated directly.
\end{proof}
This is the well-clocked RP language recorded in RP-BOOLEAN-QUOTIENT5
\cite{boolean}. Its deterministic polynomial-time status is not resolved by
the membership proof or by the geometry of this verifier's accepting tapes.

\begin{proposition}[Naive clause-product transfer fails: PROVED/EXACT]
For a nonempty proper clause $C$, let $v_l=x_i$ for a positive literal and
$v_l=1-x_i$ for a negative literal. Set
\[
 c_C=1-\prod_{l\in C}(1-v_l),\qquad P_F=\prod_{C\in F}c_C.
\]
Every CNF without empty clauses has formally nonzero $P_F$, even if UNSAT.
On Boolean assignments, however, $P_F$ equals its satisfiability indicator.
\end{proposition}
\begin{proof}
Each clause polynomial is nonzero, and a polynomial ring over a field is an
integral domain, so their product is nonzero. On Boolean points, $c_C$ is
exactly the $0/1$ truth value of its clause; their product is the conjunction.
\end{proof}
For 3-CNF with $m$ clauses, $\deg P_F\le3m$. Over a prime field with
$q\ge6m$, (4) therefore makes even an UNSAT formula evaluate nonzero with
probability at least $1/2$, violating RP NO soundness. Boolean sampling avoids
that false positive but accepts a unique-witness formula with probability
$2^{-n}$, violating YES density. The two sampling domains fail for different
reasons. The source's finite arithmetic checks remain \textbf{COMPUTATIONAL},
not a substitute for either argument.

\newpage
\section{Open links, rejected promotions, and source map}
\textbf{OPEN: language-level lower bound.} Produce one uniform $L\in\RP$
whose fixed worst-case polynomial verifier has exact one-sided soundness,
and exclude \emph{every} deterministic polynomial-time decider. A theorem
for a common clock or bounded interface remains restricted to that clock or
interface; diagonalizing over unbounded exponents must not destroy the
verifier's own fixed polynomial clock.

\textbf{OPEN: uniform simulation.} Construct a polynomial-time deterministic
decider for every RP language. If the route uses hitting lists, charge their
construction, oracle access, randomness and replay. Class membership permits
machine-dependent polynomial exponents; it does not assert a single common
exponent across every machine.

\textbf{OPEN: wrapper-to-language transfer.} A solver-independent instance or
fully charged reduction must connect the causal game to arbitrary deciders
on the \emph{same} static problem and budget. Compiling a wrapper computation
into SAT internalizes its constraints, but $\exists s,z\,G(F,s,z)$ chooses
one route. It does not impose an all-route shared representation. The paired
prime-lineage toy supplies genuine formula-internal nonmergeability, not an
arbitrary-decider representation lower bound \cite{wrapper}.

\textbf{FALSIFIED/RETIRED as proof routes:} universal public lists;
rare-witness sampling as RP; naive clause-product arithmetization; advice
existence or chosen-verifier hitting bounds as a resolution; prime/hash/timer
unpredictability as unrestricted time hardness; and finite checks as class
separations. No RP diagnostic here supplies the independent ordinary-MKTP or
fully charged ARC interfaces required by the P-versus-NP routes.

\vspace{-.5em}
\begin{thebibliography}{99}\small\setlength{\itemsep}{1pt}\setlength{\parskip}{0pt}
\bibitem{r36}\repo{research/main/2026-10-21_0848_MAIN_RP_HITTING_QUANTIFIER_FIREWALL36.tex}{RP-HITTING-QUANTIFIER-FIREWALL36}, 21 Oct. 2026.
\bibitem{r37}\repo{research/main/2026-10-21_0948_MAIN_RP_RELATIVIZATION_BIDIRECTIONAL_BARRIER37.tex}{RP-RELATIVIZATION-BIDIRECTIONAL-BARRIER37}, 21 Oct. 2026.
\bibitem{r38}\repo{research/main/2026-10-21_1048_MAIN_RP_DENSITY_HITTING_TRADEOFF38.md}{RP-DENSITY-HITTING-TRADEOFF38}, 21 Oct. 2026.
\bibitem{r39}\repo{research/main/2026-10-21_1148_MAIN_RP_LINEAR_HITTING_TIGHTNESS39.md}{RP-LINEAR-HITTING-TIGHTNESS39}, 21 Oct. 2026.
\bibitem{seed}\repo{research/review/outbox/rp-separation/2026-10-21_0812_RP_SEED_CLOCK1_status.md}{RP-SEED-CLOCK1}, 21 Oct. 2026.
\bibitem{clockfix}\repo{research/agents/lspv-1/2026-10-16_1812Z_LSPV-1_CLOCKFIX1.md}{CLOCKFIX1}, 16 Oct. 2026.
\bibitem{crs}\repo{research/explorations/random-challenge/2026-10-01_CRS-PPC-SAT1.md}{CRS-PPC-SAT1}, 1 Oct. 2026.
\bibitem{completion}\repo{research/agents/lspv-4/2026-10-14_0717_LSPV-4_COMPLETIONCUT1.tex}{COMPLETIONCUT1}, 14 Oct. 2026.
\bibitem{boolean}\repo{research/review/outbox/rp-separation/2026-10-21_1112_RP_BOOLEAN_QUOTIENT5_RP_language.md}{RP-BOOLEAN-QUOTIENT5: explicit language}; \repo{research/review/outbox/rp-separation/2026-10-21_1112_RP_BOOLEAN_QUOTIENT5_theorem.md}{clause-product theorem}, 21 Oct. 2026.
\bibitem{wrapper}\repo{research/agents/lead/2026-10-06_1436_LEAD-0023_WRAPPER_AS_SAT_QUANTIFIER_AUDIT.md}{WRAPPER-AS-SAT quantifier audit}; \repo{research/agents/lead/2026-10-06_1456_LEAD-0024_WRAPPER_AS_SAT_PRIME_LINEAGE_TOY.md}{paired prime-lineage toy}, 6 Oct. 2026.
\bibitem{gill}J. Gill. \emph{Computational Complexity of Probabilistic Turing Machines}. SICOMP 6(4), 675--695, 1977. \href{https://doi.org/10.1137/0206049}{doi:10.1137/0206049}.
\bibitem{adleman}L. Adleman. \emph{Two Theorems on Random Polynomial Time}. FOCS 1978, 75--83, Theorem I. \href{https://doi.org/10.1109/SFCS.1978.37}{doi:10.1109/SFCS.1978.37}.
\bibitem{bgs}T. Baker, J. Gill, R. Solovay. \emph{Relativizations of the P=?NP Question}. SICOMP 4(4), 431--442, 1975. \href{https://doi.org/10.1137/0204037}{doi:10.1137/0204037}.
\bibitem{bg}C. Bennett, J. Gill. \emph{Relative to a Random Oracle A, P$^A\ne$NP$^A\ne$co-NP$^A$ with Probability 1}. SICOMP 10(1), 96--113, 1981. \href{https://epubs.siam.org/doi/10.1137/0210008}{doi:10.1137/0210008}.
\bibitem{schwartz}J. T. Schwartz. \emph{Fast Probabilistic Algorithms for Verification of Polynomial Identities}. JACM 27(4), 701--717, 1980. \href{https://doi.org/10.1145/322217.322225}{doi:10.1145/322217.322225}.
\bibitem{moshkovitz}D. Moshkovitz. \emph{An Alternative Proof of the Schwartz--Zippel Lemma}. \href{https://eccc.weizmann.ac.il/report/2010/096/}{ECCC TR10-096}, 2010.
\end{thebibliography}
\end{document}
% END_OF_VERIFIED_SUMMARY::MAIN14::RP-DIAGNOSTICS
